\chapter{第十六讲\quad 刑罚的执行和消灭}
\label{lecture:16}
\noindent{\small\hyperref[ref:lecture:16]{查阅本讲扩展规则}}\par

\section{第一节\quad 减刑}
\paragraph{三个检验门}
把第78至80条拆为资格、实体、程序。资格是正在执行管制、拘役、有期徒刑或无期徒刑；死缓二年期满的转换先按第50条。实体上，认真遵守监规、接受教育改造、确有悔改表现三者合取，或者有立功表现，可以减刑；有法定重大立功，应当减刑。重大立功如舍己救人、阻止他人重大犯罪活动等，须有事实证明。程序上，执行机关向中级以上人民法院提出建议，由法院组成合议庭审理并依法裁定，三类条件分别回答“能否进入”“减刑是可还是应”“由谁依法实现”。

令 \(m\) 为减刑后累计实际执行期限的法定下界，\(L\) 为原判有限期刑期。第78条规定
\[
\begin{array}{l|l}
\text{原刑或转换情形}&m\\ \hline
\text{管制、拘役、有期徒刑}&L/2\\
\text{无期徒刑}&13\text{年}\\
\text{限制减刑的死缓期满减为无期徒刑}&25\text{年}\\
\text{限制减刑的死缓期满减为二十五年有期徒刑}&20\text{年}
\end{array}
\]
限减死缓的后两行是死缓期满后相应的最低执行期限。最高人民法院《关于办理减刑、假释案件具体应用法律的规定》第12条规定，普通死缓罪犯经减刑后实际执行不少于十五年，死缓期间不计入；第15条规定依法终身监禁者不得再减刑或假释；第18条规定缓刑一般不减刑，考验期内有重大立功者可以参照第78条减刑，同时依法缩减考验期。第80条另定无期徒刑减为有期徒刑的刑期从裁定减刑之日起计算，不能把这一新刑期起算日与累计实际执行期限的计算混作一个量。

\paragraph{例题推导}
甲原判有期徒刑八年，已依法执行两年，执行中舍己救人经查证成立。先由“有期徒刑在执行”进入第78条，再由重大立功推出应当减刑，最后算 \(m=8/2=4\) 年。假设方案使累计实际执行仅三年，则 \(3<4\)，该方案越过法定下界，不能以“应当减刑”修复；仍须经第79条程序选择合法幅度。一般悔改替换重大立功时，实体结论从“应当”降为“可以”。

\section{第二节\quad 假释}
\paragraph{先排除，再核验期限与危险}
第81条的类型门是有期徒刑或无期徒刑。对累犯，以及因故意杀人、强奸、抢劫、绑架、放火、爆炸、投放危险物质或有组织的暴力性犯罪被判十年以上有期徒刑或无期徒刑者，依法不得假释。令 \(E\) 为已实际执行期限，\(L\) 为有期徒刑原判刑期，则普通期限门为
\[
 (\text{有期}\land E\ge L/2)\ \lor\
 (\text{无期}\land E\ge 13\text{年}).
\]
还须认真遵守监规、接受教育改造、确有悔改表现、无再犯罪危险，并依法考虑社区影响及履行第82条程序。有特殊情况经最高人民法院核准，可以免除上述期限限制；该例外只指期限，不能消除禁限身份。最高人民法院上述《规定》第23、25条要求再核对死缓转换情形：非禁止对象的实际执行门槛通常为十五年，从死缓期满日起算；命中禁止条件的死缓罪犯即使减为无期或有期，仍不得假释。第83条规定考验期从假释日起算，有期徒刑的考验期为未执行完毕的刑期，无期徒刑为十年。

\paragraph{例题推导与撤销分流}
甲因题设未列入禁止范围的犯罪被判有期徒刑十二年，非累犯，实际已执行六年，符合改造和风险条件。由 \(6\ge 12/2\) 通过期限门，并依法审查社区影响，故有可以假释的入口。若同样刑期改为因抢劫罪判处，十二年达到十年门槛，命中禁止分支，已服六年不能消除禁止。若题设再给定甲依法获准假释且未减刑，考验期为 \(12-6=6\) 年。期满无第86条情形，视为原判刑罚执行完毕（第85条）；考验期内新罪撤销假释按第71条，发现漏罪撤销后按第70条，依法因违规撤销则收监执行余刑。与缓刑对读：缓刑在宣告时决定，顺利期满是原判刑罚不再执行；假释以部分执行为前提，顺利期满视为已经执行完毕。

\section{第三节\quad 刑罚的消灭与追诉时效}
\paragraph{从适用法定刑档提取时钟}
刑罚权的消灭须有法律依据，本节集中计算追诉时效。设 \(M\) 为本案依法适用的法定量刑幅度的最高刑，不用实际拟判的宣告刑代入。第87条的期限 \(T(M)\) 为
\[
 T(M)=
 \begin{cases}
 5\text{年},&M\text{ 为不满五年有期徒刑或更轻},\\
 10\text{年},&M\text{ 为五年以上且不满十年有期徒刑},\\
 15\text{年},&M\text{ 为十年以上有期徒刑},\\
 20\text{年},&M\text{ 为无期徒刑或死刑}.
 \end{cases}
\]
五年和十年的边界各进入后一档。第四档超过二十年后认为必须追诉的，须报最高人民检察院核准；其他档不能借用这一通道。共犯须分别核验适用刑档与时效事实，不能把甲的事件自动记到乙的时钟上。教材按相应刑档计算，在具体案件须先把可能改变适用档的事实和法定依据核实。

\paragraph{起点、例外、重置}
第89条第一款规定初始起点 \(t_0\)：通常取犯罪之日，有连续或继续状态则取行为终了日。第88条两条不限期路径分别是：法定机关已经立案侦查或法院受理后行为人逃避侦查或审判；被害人在期限内提出控告而法定机关应立案却未立案。单有立案这一事实未证明第一路径的逃避。依第89条第二款，在前罪期限内又犯罪时，前罪起点替换为后罪之日，期限长度仍用前罪的 \(T(M)\)。处理顺序是先检第88条，再按第89条更新普通时钟；评价日及边界须按适用法与程序事实确定。

\paragraph{例题推导}
甲第0年完成普通抢劫，题设本案法定最高刑为十年有期徒刑，故 \(T_A=15\) 年。第8年又犯普通盗窃，题设其适用法定最高刑三年，\(T_B=5\) 年，且无第88条情形。因为 \(8<15\)，前罪新起点为第8年，前罪的期限标记变为 \(8+15=23\) 年；后罪期限标记为 \(8+5=13\) 年。若第16年判断且无其他影响时效的程序事实，前罪尚在期限内，后罪已超过期限。改题为第16年才初次犯后罪，因 \(16>15\)，第89条的期限内条件失败，前罪不能因此复活。具体日期应按法定期间计算，以上整数年仅用于比较分支。

\section*{记忆与闭卷复述}
记“减看底，释看门，时效看档和钟”。闭卷先写减刑通常底线 \(1/2,13\)，再写死缓普通十五年与限减 \(25,20\) 年；假释写“刑种、禁止、半或十三、改造无险、程序、考验”；追诉写“\(5,10,15,20\) 年，先档后钟，立案加逃避，控告应立而未立，期内再犯才重开”。把“减刑实际执行下限”和“假释前已执行门槛”分别写在两列，防止数字相同导致规则错用。
\begin{enumerate}
 \item \textbf{自测。}原判有期徒刑十年，已执行四年，确有悔改且无再犯危险。普通假释期限门通过了吗？重大立功减刑是否因此一律禁止？\quad\textbf{核对。}四小于五，普通假释期限门未通过；减刑是另一规则，仍查法定事由、程序和减刑后实际执行的下限。
 \item \textbf{自测。}前案早已立案，卷内没有材料证明逃避，也没有第88条第二款情形。能否仅凭立案宣告不限期？\quad\textbf{核对。}不能。立案之后逃避是第一款合取中的另一项，要有证明。
 \item \textbf{自测。}原判缓刑与假释各自顺利期满，对原判刑罚的结论分别是什么？\quad\textbf{核对。}第76条原判刑罚不再执行；第85条原判刑罚视为已经执行完毕。
\end{enumerate}
